\subsection{Morfismos étale e inmersiones abiertas}
\label{subsection:IV.17.9}
\label{IV.17.9}

\begin{theorem}[17.9.1]
\label{IV.17.9.1}
Sea $f:X\to Y$ un morfismo. Las propiedades siguientes son equivalentes:
\begin{enumerate}
  \item[a)] $f$ es una inmersión abierta.
  \item[b)] $f$ es un monomorfismo plano localmente de presentación finita.
  \item[c)] $f$ es étale y radicial.
\end{enumerate}
\end{theorem}

Resulta de (1.4.3, (i)) que a) implica b). La condición b) implica que, para todo $y\in Y$, la fibra $f^{-1}(y)$ es vacía o isomorfa a $\Spec(k(y))$ (8.11.5.1), de modo que b) implica c), en virtud de (17.6.2, $c$)). Queda por ver que c) implica a).

Como la cuestión es local sobre $X$ y sobre $Y$ (puesto que $f$ es inyectivo), podemos limitarnos al caso en que $Y$ es afín y $f$ de presentación finita. Como $f$ es plano, es un morfismo \emph{abierto} (2.4.6), por lo que, reemplazando $Y$ por $f(X)$, podemos suponer que $f$ es \emph{sobreyectivo}. Para todo morfismo $Y'\to Y$, $f'=f_{(Y')}:X_{(Y')}\to Y'$ sigue siendo étale, radicial, sobreyectivo y de presentación finita, y por tanto abierto, y en consecuencia un homeomorfismo; dicho de otro modo, $f$ es un \emph{homeomorfismo universal}, y, al ser de tipo finito y separado (1.8.7.1), $f$ es \emph{propio}. La hipótesis de que $f$ es radicial y de tipo finito implica entonces que $f$ es cuasifinito; por tanto, (8.11.1) $f$ es un morfismo \emph{finito}. Para probar que $f$ es un isomorfismo, podemos limitarnos al caso

\oldpage[IV-4]{80}

en que $Y=\Spec(A)$, siendo $A$ un anillo local. Como $f$ es de presentación finita, se tiene $X=\Spec(B)$, donde $B$ es un $A$-módulo plano de presentación finita (1.4.7), y por tanto \emph{libre} (Bourbaki, \emph{Alg. comm.}, cap. II, § 5, n\textsuperscript{o} 2, cor. 2 del teorema 1). Además, si $\mathfrak m$ es el ideal maximal de $A$ y $k$ su cuerpo residual, $B/\mathfrak mB$ es por hipótesis un cuerpo, a la vez extensión radicial y extensión finita separable de $k$, puesto que $f$ es étale y radicial (17.6.1); por tanto, $B/\mathfrak mB$ es isomorfo a $k$, y como $B$ es un $A$-módulo libre, $B$ es isomorfo a $A$. C.Q.D.

\begin{corollary}[17.9.2]
\label{IV.17.9.2}
Sea $X$ un preesquema conexo; si $f:X\to Y$ es una inmersión cerrada étale, entonces $f$ es un isomorfismo de $X$ sobre una componente conexa abierta de $Y$.
\end{corollary}

En efecto, $f$ es étale y radicial, y por tanto un isomorfismo de $X$ sobre un preesquema inducido en un subconjunto abierto de $Y$; pero, por hipótesis, $f(X)$ es cerrado en $Y$, y por tanto a la vez abierto y cerrado, y como $f(X)$ es conexo, es una componente conexa de $Y$.

\begin{corollary}[17.9.3]
\label{IV.17.9.3}
Sea $f:X\to Y$ un morfismo étale (resp. étale y separado). Entonces toda $Y$-sección $g:Y\to X$ de $X$ es una inmersión abierta (resp. abierta y cerrada); además, la aplicación $g\mapsto g(Y)$ es una biyección del conjunto $\Gamma(X/Y)$ de las $Y$-secciones de $X$ sobre el conjunto de los subconjuntos abiertos $Z$ (resp. abiertos y cerrados) de $X$ tales que la restricción de $f$ a $Z$ sea un morfismo sobreyectivo y radicial de $Z$ sobre $Y$.
\end{corollary}

En efecto, el hecho de que $g$ sea una inmersión abierta resulta ya de que $f$ es no ramificado (17.4.1, $b''$)), y la restricción de $f$ al abierto $g(Y)$ de $X$ es un isomorfismo. Recíprocamente, si $Z$ es un abierto de $X$ tal que $f|Z$ sea un morfismo sobreyectivo y radicial de $Z$ sobre $Y$, $f|Z$ es un isomorfismo en virtud de (17.9.1), puesto que es étale. Si $f$ es además separado, se sabe que $g$ es una inmersión cerrada (1, 5.4.6), lo que termina de probar el corolario.

\begin{corollary}[17.9.4]
\label{IV.17.9.4}
Sean $Y$ un preesquema conexo y $f:X\to Y$ un morfismo étale y separado. Entonces toda $Y$-sección $g$ de $X$ es un isomorfismo de $Y$ sobre una componente conexa abierta de $X$, y la aplicación $g\mapsto g(Y)$ es una biyección de $\Gamma(X/Y)$ sobre el conjunto de las componentes conexas abiertas $Z$ de $X$ tales que la restricción de $f$ a $Z$ sea un morfismo sobreyectivo y radicial de $Z$ sobre $Y$.
\end{corollary}

\begin{corollary}[17.9.5]
\label{IV.17.9.5}
Sean $g:Y\to S$, $h:X\to S$ dos morfismos localmente de presentación finita; se supone además que $h$ es plano. Para que un $S$-morfismo $f:X\to Y$ sea una inmersión abierta (resp. un isomorfismo), es necesario y suficiente que, para todo $s\in S$, el morfismo $f_s:h^{-1}(s)\to g^{-1}(s)$ deducido de $f$ por el cambio de base $\Spec(k(s))\to S$ sea una inmersión abierta (resp. un isomorfismo).
\end{corollary}

En efecto, si $f_s$ es una inmersión abierta para todo $s\in S$, resulta de (17.8.3) que $f$ es un morfismo étale; como, para todo $y\in Y$, se tiene $f^{-1}(y)=f_s^{-1}(y)$ con $s=g(y)$, $f$ es radicial; por tanto, $f$ es una inmersión abierta en virtud de (17.9.1). Si, además, $f_s$ es sobreyectivo para todo $s\in S$, $f$ es sobreyectivo, y por tanto un isomorfismo.

La proposición siguiente precisa (10.4.11):

\begin{proposition}[17.9.6]
\label{IV.17.9.6}
Sean $S$ un preesquema y $X$ un $S$-preesquema de presentación finita. Todo $S$-endomorfismo de $X$ que sea un monomorfismo es un automorfismo de $X$.
\end{proposition}

Sean $f:X\to S$ el morfismo estructural y $g$ el $S$-endomorfismo considerado. La cuestión es local sobre $S$, por lo que podemos suponer que $S$ es afín. Utilizando (8.9.1) y (8.10.5, (i \emph{bis})), nos reducimos al caso en que $S=\Spec(A)$, donde $A$ es una $\mathbb{Z}$-álgebra de tipo finito, y por tanto $X$ es un $\mathbb{Z}$-preesquema de tipo finito. Resulta ya de (10.4.11) que $g$ es un morfismo biyectivo, puesto que un monomorfismo es radicial (8.11.5.1); bastará, por tanto, mostrar

\oldpage[IV-4]{81}

que $g$ es una inmersión abierta, y, puesto que $g$ es radicial, bastará, en virtud de (17.9.1), probar que $g$ es étale. Además, como el conjunto de los puntos donde $g$ es étale es abierto y $X$ es un preesquema de Jacobson (10.4.7), bastará mostrar que $g$ es étale en todo punto \emph{cerrado} $z$ de $X$ (10.3.1). Pongamos $z'=g(z)$; resulta de (10.4.11.1, (i)) que $z'$ también es un punto cerrado de $X$ y, puesto que $g$ es un monomorfismo, la aplicación $k(z')\to k(z)$ deducida de $g$ es un \emph{isomorfismo}. Para que $g$ sea étale en el punto $z$, es necesario y suficiente, por tanto, que el homomorfismo canónico $\widehat{\sh{O}}_{X,z'}\to\widehat{\sh{O}}_{X,z}$ sea biyectivo (17.6.3, $e''$)). Para ello, probaremos que, para todo entero $n$, el homomorfismo canónico $\sh{O}_{X,z'}/\mathfrak m_{z'}^{n+1}\to\sh{O}_{X,z}/\mathfrak m_z^{n+1}$ es biyectivo, de donde resultará inmediatamente la conclusión. Podemos suponer que $z$ pertenece al conjunto finito $T_{p,d}$ de los puntos cerrados $t$ de $X$ tales que $k(t)$ sea una extensión de $\mathbb{F}_p$ cuyo grado divide a $d$; en cuyo caso lo mismo ocurre con $z'$. Designemos por $\sh{J}$ el ideal coherente de $\sh{O}_X$ tal que $\sh{J}_x=\sh{O}_x$ para $x\notin T_{p,d}$, y $\sh{J}_x=\mathfrak m_x^{n+1}$ para $x\in T_{p,d}$. Sean $T_{p,d,n}$ el subpreesquema cerrado de $X$ definido por $\sh{J}$, y $j:T_{p,d,n}\to X$ la inyección canónica; el morfismo compuesto $g\circ j:T_{p,d,n}\to X$ lleva el conjunto $T_{p,d}$ a sí mismo, y para todo $x\in T_{p,d}$, el homomorfismo $\sh{O}_{X,g(x)}\to\sh{O}_{X,x}/\mathfrak m_x^{n+1}$ se factoriza como
\[
\sh{O}_{X,g(x)}\to\sh{O}_{X,g(x)}/\mathfrak m_{g(x)}^{n+1}\to\sh{O}_{X,x}/\mathfrak m_x^{n+1}.
\]
Por tanto (I, 4.9), existe un único endomorfismo $g'$ de $T_{p,d,n}$ tal que $j\circ g'=g\circ j$. Como $j$ es un monomorfismo, y lo mismo ocurre con $g$ por hipótesis, se deduce que $g'$ también es un monomorfismo, y la proposición quedará probada si se muestra que $g'$ es un \emph{automorfismo} de $T_{p,d,n}$. Ahora bien, para todo $x\in T_{p,d,n}$, hemos visto que $k(x)$ es un cuerpo finito, y como $\sh{O}_{X,x}$ es noetheriano, cada uno de los $\mathfrak m_x^h/\mathfrak m_x^{h+1}$ es un $k(x)$-espacio vectorial de dimensión finita; de donde se concluye inmediatamente que el anillo local $\sh{O}_{X,x}/\mathfrak m_x^{n+1}$ de $T_{p,d,n}$ en el punto $x$ tiene un número \emph{finito} de elementos y, \emph{a fortiori}, es un $\mathbb{Z}$-módulo de tipo finito. Como $T_{p,d,n}$ es suma de un número finito de preesquemas $\Spec(\sh{O}_{X,x}/\mathfrak m_x^{n+1})$, es un $\mathbb{Z}$-preesquema \emph{finito}; el endomorfismo $g'$ es, por tanto, también \emph{finito} (II, 6.1.5, (v)), y por consiguiente propio. Pero un monomorfismo propio es una \emph{inmersión cerrada} (8.11.5); si $T_{p,d,n}=\Spec(B)$, $g'$ corresponde entonces a un endomorfismo sobreyectivo $\varphi:B\to B$ del anillo $B$. Ahora bien, el conjunto $B$ es finito, por lo que $\varphi$ es necesariamente biyectivo. C.Q.D.
